%! AP Statistics — Unit 4, Lesson 4.12 — Guided Notes
\documentclass[10pt]{article}
\usepackage{apstats-article}
\usepackage{apstats-boxes}

\begin{document}

\pageheader{Unit 4, Lesson 4.12}{Guided Notes}
\namedateperiod

\vspace{0.12in}

\begin{objectivebox}
\small By the end of this lesson, I will be able to\dots\\[4pt]
\writeline\\[1pt]
\writeline\\[1pt]
\writeline
\end{objectivebox}

\vspace{0.1in}

\begin{vocabbox}
\termblanklong{Geometric random variable}
\termblanklong{BITS conditions}
\termblanklong{Geometric probability formula}
\termblanklong{Mean of a geometric distribution}
\termblanklong{geometpdf(p, k) / geometcdf(p, k)}
\end{vocabbox}

\vspace{0.1in}

\begin{hookbox}
\small\textbf{Opening scenario:} A basketball player makes 65\% of free throws.
He keeps shooting until he makes one.

\vspace{6pt}
On average, how many shots do you expect him to need? \blank{3cm}\\[6pt]
What values can $X$ take? \blank{6cm} Is there an upper limit? \blank{2cm}
\end{hookbox}

\vspace{0.1in}

\begin{notesbox}{Geometric Settings --- BITS Conditions}
\small
A random variable $X$ follows a \textbf{geometric distribution} if all four
\textbf{BITS} conditions are met. Compare to BINS:

\vspace{8pt}
\renewcommand{\arraystretch}{2.0}
\begin{tabularx}{\linewidth}{@{} >{\bfseries\color{navy}}p{0.06\linewidth}
                                     >{\bfseries}p{0.30\linewidth} X
                                     p{3.2cm} @{}}
 & \textbf{BITS condition} & \textbf{What it means} & \textbf{Binomial analog} \\
\hline
B & Binary & \blank{4cm} & same \\
I & Independent & \blank{4cm} & same \\
T & Trials until \blank{2cm} success & Counting stops when the first \blank{2cm} occurs & N: fixed $n$ \\
S & Same $p$ & \blank{4cm} & same \\
\end{tabularx}

\vspace{6pt}
If all four conditions are met, write: $X \sim \text{Geom}(\blank{1.5cm})$
\end{notesbox}

\vspace{0.1in}

\begin{notesbox}{The Geometric Probability Formula and Mean}
\small

\textbf{Formula:} For $X \sim \text{Geom}(p)$, the probability that the first
success occurs on trial $k$ is:
\[
  P(X = k) = (1-p)^{\blank{1cm}} \cdot p \qquad k = 1, 2, 3, \ldots
\]

\vspace{4pt}
Why? Trials $1$ through $k-1$ must be \blank{3cm}, and trial $k$ must be a
\blank{2cm}.

\vspace{8pt}
\textbf{Mean:} $\mu_X = \dfrac{1}{\blank{1cm}}$

\vspace{4pt}
\textbf{Interpretation:} ``In many repetitions of this geometric experiment, the
average number of \blank{3cm} until the first \blank{3cm} would be $1/p$.''

\vspace{8pt}
\textbf{Calculator:}
\begin{itemize}[itemsep=4pt]
  \item \texttt{geometpdf(p, k)} $= P(X = k)$
  \item \texttt{geometcdf(p, k)} $= P(X \le k)$
  \item $P(X > k) = (1-p)^k = 1 - \texttt{geometcdf(p, k)}$
\end{itemize}
\end{notesbox}

\vspace{0.1in}

\begin{notesbox}{Worked Example --- Free-Throw Shooting}
\small
\textbf{Context:} Basketball player, $p = 0.65$, shoots until first make.\\
$X = $ \blank{7cm} \quad $X \sim \text{Geom}(\blank{1.5cm})$

\vspace{6pt}
\textbf{Check BITS:}\\[4pt]
B: \blank{10cm}\\[4pt]
I: \blank{10cm}\\[4pt]
T: \blank{10cm}\\[4pt]
S: \blank{10cm}

\vspace{8pt}
\textbf{Compute} $P(X = 3)$ (first make on the 3rd shot):
\[
  P(X = 3) = (1 - \blank{1.5cm})^{\blank{1cm}} \cdot \blank{1.5cm}
           = (\blank{1.5cm})^2 \cdot \blank{1.5cm}
           = \blank{3cm}
\]
Calculator: \texttt{geometpdf(}\blank{1cm}\texttt{,}\;3\texttt{)} $= $ \blank{2cm}

\vspace{8pt}
\textbf{Mean:} $\mu_X = \dfrac{1}{\blank{1.5cm}} \approx \blank{2cm}$

\vspace{6pt}
\textbf{Interpret $\mu_X$:} \writeline\\[1pt]\writeline
\end{notesbox}

\vspace{0.1in}

\begin{practicebox}
\small
\begin{enumerate}[leftmargin=1.5em, itemsep=10pt]

  \item A standard die is rolled until a 6 appears. Let $X =$ the number of
  rolls. Assume $X \sim \text{Geom}(1/6)$.
  \begin{enumerate}[label=(\alph*), itemsep=6pt]
    \item Find $P(X = 4)$. Show formula setup.\\
          \writeline
    \item Find $\mu_X$ and interpret in context.\\
          \writeline\\[1pt]\writeline
    \item Find $P(X > 4) = (1-p)^4$. Interpret in one sentence.\\
          \writeline\\[1pt]\writeline
  \end{enumerate}

\end{enumerate}
\end{practicebox}

\end{document}
